% !TEX root = ../../../main.tex
% 所属：第四部分：Fourier 分析 / 第5章：$\R^d$上的Fourier变换
% 节：5.1 预备知识
% 正文按原文件顺序保存；由 ../chapter5.tex 调用。
\section{预备知识}

这一章考虑的空间为 $\R^d$，它是由所有 $d$ 元数组
\[
x=(x_1,x_2,\ldots,x_d)
\]
组成的，其中 $x_i\in\R$。向量加法定义为对应分量相加，向量数乘的定义类似。对于
$x=(x_1,x_2,\ldots,x_d)\in\R^d$，定义
\[
\abs{x}=\left(x_1^2+\cdots+x_d^2\right)^{1/2},
\]
使 $\abs{x}$ 正是一般欧氏范数下向量 $x$ 的长度。实际上，赋予 $\R^d$ 的内积为
\[
x\cdot y=x_1y_1+\cdots+x_dy_d,
\]
使得 $\abs{x}^2=x\cdot x$。在$\R^d$上，我们用 $x\cdot y$ 替代$\langle x,y\rangle$。

\begin{definition}
给定 $d$ 元非负整数数组
\[
\alpha=(\alpha_1,\ldots,\alpha_d),
\]
称其为一个多重指标，单项式 $x^\alpha$ 定义为
\[
x^\alpha=x_1^{\alpha_1}x_2^{\alpha_2}\cdots x_d^{\alpha_d}.
\]
类似地，定义微分算子 $(\partial/\partial x)^\alpha$ 为
\[
\left(\frac{\partial}{\partial x}\right)^\alpha
=\left(\frac{\partial}{\partial x_1}\right)^{\alpha_1}
 \left(\frac{\partial}{\partial x_2}\right)^{\alpha_2}
 \cdots
 \left(\frac{\partial}{\partial x_d}\right)^{\alpha_d}
=\frac{\partial^{\abs{\alpha}}}
{\partial x_1^{\alpha_1}\cdots\partial x_d^{\alpha_d}},
\]
其中
\[
\abs{\alpha}=\alpha_1+\cdots+\alpha_d
\]
是多重指标 $\alpha$ 的阶数。
\end{definition}

空间 $\R^d$ 上的分析，特别是 Fourier 变换理论，由底空间的三个重要的对称群所塑造：
\begin{enumerate}[label=(\roman*)]
	\item 平移；
	\item 伸缩；
	\item 旋转。
\end{enumerate}
我们已经看到平移 $x\longmapsto x+h$，其中 $h\in\R^d$ 固定，以及伸缩
$x\longmapsto\delta x$，其中 $\delta>0$。
\newpage
\begin{definition}
	$\R^d$ 中的旋转是保持内积的线性变换
	\[
	\mathcal{R}:\R^d\longrightarrow\R^d.
	\]
	换而言之，
	\begin{enumerate}[label=\textbullet]
		\item 对所有 $x,y\in\R^d$ 和 $a,b\in\R$，有
		\[
		\mathcal{R}(ax+by)=a\mathcal{R}(x)+b\mathcal{R}(y);
		\]
		\item 对所有 $x,y\in\R^d$，有
		\[
		\mathcal{R}(x)\cdot \mathcal{R}(y)=x\cdot y.
		\]
	\end{enumerate}
\end{definition}

第二个条件等价于
\[
\abs{\mathcal{R}(x)}=\abs{x},\qquad x\in\R^d,
\]
也等价于
\[
\mathcal{R}^{ T}=\mathcal{R}^{-1},
\]
其中 $\mathcal{R}^{ T}$ 和 $\mathcal{R}^{-1}$ 分别表示 $\mathcal{R}$ 的转置和逆变换。特别地，
\[
\det(\mathcal{R})=\pm1,
\]
其中 $\det(\mathcal{R})$ 表示 $\mathcal{R}$ 的行列式。如果 $\det(\mathcal{R})=1$，则称 $\mathcal{R}$ 为正常旋转；否则称其为反常旋转。

\begin{definition}
	$\R^d$ 上的连续复值函数 $f$ 称为\textbf{速降的}，如果对于每个多重指标 $\alpha$，函数
	\[
	\abs{x^\alpha f(x)}
	\]
	是有界的。等价地，连续函数 $f$ 为速降函数，如果
	\[
	\sup_{x\in\R^d}\abs{x}^k\abs{f(x)}<\infty,
	\qquad k=0,1,2,\ldots.
	\]
\end{definition}

\newpage
\subsection*{积分}
\begin{definition}
	给定一个速降函数$f$，定义
\[
\int_{\R^d}f(x)\dd x
:=\lim_{N\to\infty}\int_{Q_N}f(x)\dd x,
\]
其中 $Q_N$ 表示中心在原点、边长为 $N$ 的闭立方体，并且立方体的各边平行于坐标轴，即
\[
Q_N=\left\{x\in\R^d:\abs{x_i}\leq\frac{N}{2},\ i=1,\ldots,d\right\}.
\]
立方体 $Q_N$ 上的积分为通常意义下 Riemann 积分。由于上述积分
\[
I_N=\int_{Q_N}f(x)\dd x
\]
当 $N$ 趋于无穷时构成 Cauchy 序列，故上述极限存在
\end{definition}

事实上，取 $k>d$。由速降性，存在常数 $C>0$，使得当 $\abs{x}\geq1$ 时，
\[
\abs{f(x)}\leq C\abs{x}^{-k}.
\]
而立方壳 $Q_{n+1}\setminus Q_n$ 的体积为 $O(n^{d-1})$，故积分的尾项可由收敛级数
$\sum_{n=1}^{\infty}n^{d-1-k}$ 控制，从而 $\{I_N\}$ 是 Cauchy 序列。

依次有下面两个观察：

首先，可以用球
\[
B_N=\left\{x\in\R^d:\abs{x}\leq N\right\}
\]
来替代立方体 $Q_N$ 而不改变定义；

其次，不需要速降函数的所有条件来保证极限的存在。实际上，只需假设 $f$ 是连续的，并且对某个 $\eps>0$ 有
\[
\sup_{x\in\R^d}\abs{x}^{d+\eps}\abs{f(x)}<\infty.
\]
例如，$\R$ 上适度下降函数对应于 $\eps=1$
\begin{definition}
	定义 $\R^d$ 上的适度下降函数为满足
\[
\sup_{x\in\R^d}\abs{x}^{d+1}\abs{f(x)}<\infty
\]
的连续函数。
\end{definition}

积分与三个重要对称群的相互作用如下。

\begin{proposition}
	若 $f$ 是 $\R^d$ 上的适度下降函数，那么：
	\begin{enumerate}[label=(\roman*)]
		\item 对任意 $h\in\R^d$，
		\[
		\int_{\R^d}f(x+h)\dd x=\int_{\R^d}f(x)\dd x;
		\]

		\item 对任意 $\delta>0$，
		\[
		\delta^d\int_{\R^d}f(\delta x)\dd x
		=\int_{\R^d}f(x)\dd x;
		\]

		\item 对任意旋转 $R$，
		\[
		\int_{\R^d}f(R(x))\dd x
		=\int_{\R^d}f(x)\dd x.
		\]
	\end{enumerate}
\end{proposition}

\begin{proof}
	三式均由多重积分的变量替换得到。分别令
	\[
	u=x+h,\qquad u=\delta x,\qquad u=R(x).
	\]
	相应的 Jacobian 行列式的绝对值依次为
	\[
	1,\qquad \delta^d,\qquad \abs{\det(R)}=1,
	\]
	代入即得结论。
\end{proof}

\subsection*{极坐标}

在 $\R^d$ 中引进极坐标并找出相应的积分公式将为我们提供便利。先以 $d=2$ 和 $d=3$ 的情形为例

\begin{example}
	在 $\R^2$ 上，极坐标由 $(r,\theta)$ 给出，其中 $r\geq0$ 且 $0\leq\theta<2\pi$。变量代换
	\[
	(r,\theta)\longmapsto(r\cos\theta,r\sin\theta)
	\]
	的 Jacobian 行列式为 $r$，因此
	\[
	\int_{\R^2}f(x)\dd x
	=\int_0^{2\pi}\int_0^\infty
	f(r\cos\theta,r\sin\theta)\,r\dd r\dd\theta.
	\]
\end{example}

现将单位圆 $S^1$ 上的点写成
\[
\gamma=(\cos\theta,\sin\theta).
\]
考虑圆 $S^1$ 上的函数 $g$，定义它在 $S^1$ 上的积分为
\[
\int_{S^1}g(\gamma)\dd\sigma(\gamma)
:=\int_0^{2\pi}g(\cos\theta,\sin\theta)\dd\theta.
\]
利用此记号，有
\[
\int_{\R^2}f(x)\dd x
=\int_{S^1}\int_0^\infty f(r\gamma)\,r\dd r\dd\sigma(\gamma).
\]

\begin{example}
	在 $\R^3$ 上，使用如下球面坐标公式：
	\[
	\begin{cases}
		x_1=r\sin\theta\cos\varphi,\\
		x_2=r\sin\theta\sin\varphi,\\
		x_3=r\cos\theta,
	\end{cases}
	\]
	其中 $r>0$，$0\leq\theta\leq\pi$ 且 $0\leq\varphi\leq2\pi$。变量代换的 Jacobian 行列式为
	$r^2\sin\theta$，因此
	\[
	\int_{\R^3}f(x)\dd x
	=\int_0^{2\pi}\int_0^\pi\int_0^\infty
	f(r\sin\theta\cos\varphi,
	  r\sin\theta\sin\varphi,
	  r\cos\theta)
	\,r^2\dd r\sin\theta\dd\theta\dd\varphi.
	\]
\end{example}

如果 $g$ 是单位球面
\[
S^2=\left\{x\in\R^3:\abs{x}=1\right\}
\]
上的函数，令
\[
\gamma=(\sin\theta\cos\varphi,
        \sin\theta\sin\varphi,
        \cos\theta),
\]
定义球面积分为
\[
\int_{S^2}g(\gamma)\dd\sigma(\gamma)
:=\int_0^{2\pi}\int_0^\pi
g(\gamma)\sin\theta\dd\theta\dd\varphi.
\]
因此，
\[
\int_{\R^3}f(x)\dd x
=\int_{S^2}\int_0^\infty
f(r\gamma)r^2\dd r\dd\sigma(\gamma).
\]

一般地，可将 $\R^d\setminus\{0\}$ 的任一点唯一表示为
\[
x=r\gamma,
\]
其中 $\gamma$ 在球面 $S^{d-1}\subset\R^d$ 上且 $r>0$。实际上，取
\[
r=\abs{x},\qquad \gamma=\frac{x}{\abs{x}}.
\]
如同 $d=2$ 或 $d=3$ 的情形，可以接着定义一般的球面坐标。用到的公式为
\[
\int_{\R^d}f(x)\dd x
=\int_{S^{d-1}}\int_0^\infty
f(r\gamma)r^{d-1}\dd r\dd\sigma(\gamma),
\]
只要 $f$ 是适度下降函数。此处 $\dd\sigma(\gamma)$ 表示由球面坐标导出的球面 $S^{d-1}$ 中的面积元。

